% =========================
% chapters/08_discussion_future_conclusion.tex
% =========================
\chapter{Discussion, Challenges, Future Work, and Conclusion}
\label{ch:discussion_future}

\section{Discussion: what Phaser enables in practice}
This thesis argues that Phaser’s architecture (DAC-based AWG platform with flexible digital mixing and
gateware-defined modes) is well-matched to modern trapped-ion control demands such as shaped pulses and
multitone waveforms, provided synchronization and timing semantics are handled carefully.

% TODO: Add a short list of “worked well” items and connect to validation plots.

\section{Challenges and lessons learned}
This section should document practical realities:
(i) initialization order and “hidden state” issues,
(ii) sync triggers and update staging (why a config wasn’t applied when expected),
(iii) calibration needs (IQ imbalance, cable delays),
and (iv) software engineering constraints (kernel restrictions, host/kernel data movement).

% TODO: convert your debugging narrative into 2–3 pages of crisp engineering lessons.

\section{Future work}
Future work should be prioritized by impact:
1) automated calibration and compensation (IQ, amplitude flatness, phase offsets),
2) closed-loop control using onboard ADC channels,
3) scaling experiments: more channels, satellites, or DRTIO scenarios,
4) more formal verification of timing/phase claims (continuous integration with waveform simulation).

\section{Conclusion}
Summarize (i) the system you built, (ii) what was validated, and (iii) how it improves experimental capability.
